\documentclass[11pt,a4paper]{article}
% Même mise en forme que le polycopié de prérequis (boîtes, exercices, corrigés).
% ===========================================================================
%  Polycopié de prérequis — Comptabilité prudentielle (M2 alternance)
%  Préambule partagé
%  F. Borel-Mathurin — Université Paris 1 Panthéon-Sorbonne
% ===========================================================================

\usepackage[utf8]{inputenc}
\usepackage[T1]{fontenc}
\IfFileExists{french.ldf}{\usepackage[french]{babel}}{}
\IfFileExists{lmodern.sty}{\usepackage{lmodern}}{}
\usepackage[a4paper,margin=2.4cm]{geometry}

\usepackage{amsmath,amssymb}
\usepackage{booktabs}
\usepackage{array}
\usepackage{enumitem}
\usepackage{xcolor}
\usepackage{tcolorbox}
\tcbuselibrary{skins,breakable}
\usepackage{tikz}
\usetikzlibrary{arrows.meta,positioning,calc}
\usepackage{fancyhdr}
\usepackage[hidelinks]{hyperref}
\usepackage{makeidx}
\makeindex
\IfFileExists{eurosym.sty}{\usepackage{eurosym}}{\providecommand{\euro}{\texteuro}}

% ---------------------------------------------------------------- couleurs
% À remplacer par les valeurs exactes de la charte si besoin.
\definecolor{bleuactu}{RGB}{0,84,150}
\definecolor{vertactu}{RGB}{0,132,100}
\definecolor{rougeactu}{RGB}{175,30,45}
\definecolor{orangeactu}{RGB}{230,140,0}
\definecolor{grisactu}{RGB}{110,110,110}

% ------------------------------------------------- version étudiant/corrigé
% \correctiontrue  -> corrigés inclus
% \correctionfalse -> version étudiant
\newif\ifcorrection
\correctiontrue

% ------------------------------------------------------------------ boîtes
\newtcolorbox{notion}[1]{%
  breakable, enhanced, colback=bleuactu!4, colframe=bleuactu,
  boxrule=0.4pt, left=8pt, right=8pt, top=6pt, bottom=6pt,
  fonttitle=\bfseries\sffamily, coltitle=white, title={#1}}

\newtcolorbox{generalisation}[1]{%
  breakable, enhanced, colback=orangeactu!5, colframe=orangeactu!85!black,
  boxrule=0.4pt, left=8pt, right=8pt, top=6pt, bottom=6pt,
  fonttitle=\bfseries\sffamily, coltitle=white, title={#1}}

\newtcolorbox{pont}[1]{%
  breakable, enhanced, colback=rougeactu!4, colframe=rougeactu,
  boxrule=0.4pt, left=8pt, right=8pt, top=6pt, bottom=6pt,
  fonttitle=\bfseries\sffamily, coltitle=white, title={#1}}

\newtcolorbox{refreg}{%
  breakable, enhanced, colback=grisactu!6, colframe=grisactu!60,
  boxrule=0.4pt, left=8pt, right=8pt, top=6pt, bottom=6pt,
  fontupper=\small}

% --------------------------------------------------------------- exercices
\newcounter{exo}[section]
\renewcommand{\theexo}{\thesection.\arabic{exo}}

\newenvironment{exercice}[1][]{%
  \refstepcounter{exo}%
  \par\medskip\noindent
  \begin{tcolorbox}[breakable, enhanced, colback=white, colframe=vertactu,
     boxrule=0.4pt, left=8pt, right=8pt, top=6pt, bottom=6pt,
     fonttitle=\bfseries\sffamily, coltitle=white,
     title={Exercice \theexo\ifx\relax#1\relax\else\ --- #1\fi}]}%
  {\end{tcolorbox}\par\medskip}

\newenvironment{corrige}{%
  \ifcorrection
    \par\smallskip\noindent
    \begin{tcolorbox}[breakable, enhanced, colback=vertactu!4,
       colframe=vertactu!50, boxrule=0.4pt, left=8pt, right=8pt,
       top=6pt, bottom=6pt, fontupper=\small,
       title={\itshape Corrigé}, fonttitle=\small\itshape,
       coltitle=vertactu!80!black, colbacktitle=vertactu!10]%
  \else
    \setbox0\vbox\bgroup
  \fi}%
  {\ifcorrection\end{tcolorbox}\par\smallskip\else\egroup\fi}

% ------------------------------------------------------------- échéanciers
\tikzset{
  timeaxis/.style={-{Latex[length=2mm]}, thick, grisactu},
  timenode/.style={circle, fill=bleuactu, inner sep=1.1pt},
  flux/.style={-{Latex[length=2mm]}, thick, vertactu},
  fluxneg/.style={-{Latex[length=2mm]}, thick, rougeactu},
  etiq/.style={font=\scriptsize, inner sep=1pt}
}

% ------------------------------------------------------------------ titres
\usepackage{titlesec}
\titleformat{\section}{\Large\bfseries\sffamily\color{bleuactu}}{\thesection}{0.7em}{}
\titleformat{\subsection}{\large\bfseries\sffamily\color{bleuactu!80!black}}{\thesubsection}{0.6em}{}

\pagestyle{fancy}
\fancyhf{}
\fancyhead[L]{\small\sffamily\color{grisactu}Prérequis --- Comptabilité prudentielle}
\fancyhead[R]{\small\sffamily\color{grisactu}\thepage}
\renewcommand{\headrulewidth}{0.2pt}

\newcommand{\notation}[2]{\item[{$#1$}] #2}  % accolades : #1 peut contenir des crochets


% Version étudiant par défaut ; `make pdf` produit aussi td-corrige.pdf en définissant \Corrige.
\ifdefined\Corrige\correctiontrue\else\correctionfalse\fi

\fancyhead[L]{\small\sffamily\color{grisactu}Séance 1 --- Comptabilité prudentielle}
\newcommand{\keur}{\,k\euro}

\title{\vspace{-1.5cm}\sffamily\bfseries\color{bleuactu}
       TD 1 : le mini-bilan prudentiel de l'Assureur Bêta\\[2mm]
       \large\mdseries\color{black}Du bilan des comptes sociaux au ratio de couverture du SCR
       \ifcorrection\\[1mm]\color{vertactu}Version corrigée\fi}
\author{\sffamily Fabrice Borel-Mathurin}
\date{\sffamily\small Master SAFIR --- 2026-2027}

\begin{document}
\maketitle
\thispagestyle{empty}

\begin{notion}{Objectifs du TD (2 heures, sur tableur)}
À partir du bilan des comptes sociaux d'un assureur vie fictif, construire son bilan prudentiel,
classer ses fonds propres, vérifier la couverture du SCR et du MCR, puis mesurer la sensibilité du
résultat à trois chocs de marché. Les données sont dans \texttt{assureur-beta.csv} (séparateur
« ; », s'ouvre directement dans Excel). Montants en milliers d'euros (k\euro).
\end{notion}

\section{Données}

\begin{center}
\small
\begin{tabular}{lrr@{\qquad}lr}
\toprule
\textbf{Actif} & \textbf{Comptable} & \textbf{Marché} & \textbf{Passif} & \textbf{Comptable} \\
\midrule
Obligations d'État        & 500\,000 & 480\,000 & Provisions techniques       & 950\,000 \\
Obligations d'entreprises & 200\,000 & 195\,000 & Dettes subordonnées (12 ans) & 50\,000 \\
Actions cotées            & 100\,000 &  90\,000 & Autres dettes               &  20\,000 \\
Immobilier                & 200\,000 & 210\,000 & Capital social              &  30\,000 \\
Trésorerie                &  50\,000 &  50\,000 & Réserves et résultat        &  20\,000 \\
Créances                  &  20\,000 &  19\,000 &                             &          \\
\midrule
\textbf{Total}            & \textbf{1\,070\,000} & & \textbf{Total} & \textbf{1\,070\,000} \\
\bottomrule
\end{tabular}
\end{center}

\noindent Hypothèses :
\begin{itemize}[nosep]
  \item la meilleure estimation des engagements vaut 920\,000\keur{} ;
  \item la base fiscale de chaque poste est sa valeur dans les comptes sociaux ; taux d'impôt 25\,\% ;
        les impôts différés actifs sont recouvrables ;
  \item la dette subordonnée remplit les critères du niveau 2 et reste valorisée à 50\,000\keur{} ;
  \item SCR = 70\,000\keur{} et MCR = 17\,500\keur{} (leur calcul est l'objet de la séance 7).
\end{itemize}

\begin{exercice}[Revalorisation des actifs]
\begin{enumerate}[nosep]
  \item Calculer, poste par poste, l'écart entre valeur de marché et valeur comptable, puis l'écart total.
  \item Proposer une explication économique à chaque écart non nul.
\end{enumerate}
\end{exercice}
\begin{corrige}
\begin{center}
\begin{tabular}{lrrr}
\toprule
Poste & Comptable & Prudentiel & Écart \\
\midrule
Obligations d'État        & 500\,000 & 480\,000 & $-20\,000$ \\
Obligations d'entreprises & 200\,000 & 195\,000 & $-5\,000$ \\
Actions cotées            & 100\,000 &  90\,000 & $-10\,000$ \\
Immobilier                & 200\,000 & 210\,000 & $+10\,000$ \\
Trésorerie                &  50\,000 &  50\,000 & 0 \\
Créances                  &  20\,000 &  19\,000 & $-1\,000$ \\
\midrule
Total                     & 1\,070\,000 & 1\,044\,000 & $\mathbf{-26\,000}$ \\
\bottomrule
\end{tabular}
\end{center}
Obligations : moins-values latentes après une hausse des taux (État) et un élargissement des
écarts de crédit (entreprises). Actions : baisse des marchés. Immobilier : plus-value latente sur
un actif inscrit au coût historique. Créances : risque de non-recouvrement. L'article~75 de la
directive impose la valeur d'échange ; la hiérarchie de valorisation (règlement délégué, art.~10)
part du prix coté sur un marché actif.
\end{corrige}

\begin{exercice}[Provisions techniques]
\begin{enumerate}[nosep]
  \item Avec un coût du capital de 6\,\%, estimer la marge de risque à partir de la projection
        simplifiée du SCR ci-dessous (marge de la tranche = coût du capital $\times$ nombre
        d'années $\times$ SCR moyen $\times$ facteur d'actualisation moyen). Arrondir au millier.
  \item En déduire les provisions techniques prudentielles et leur écart avec les provisions comptables.
\end{enumerate}
\begin{center}
\small
\begin{tabular}{lrrr}
\toprule
Années & Nombre d'années & SCR moyen projeté & Facteur d'actualisation moyen \\
\midrule
1 à 5   & 5  & 70\,000 & 0,92 \\
6 à 10  & 5  & 50\,000 & 0,75 \\
11 à 20 & 10 & 30\,000 & 0,55 \\
\bottomrule
\end{tabular}
\end{center}
\end{exercice}
\begin{corrige}
$\mathrm{MR} = 6\,\% \times (5 \times 70\,000 \times 0{,}92 + 5 \times 50\,000 \times 0{,}75
 + 10 \times 30\,000 \times 0{,}55) = 19\,320 + 11\,250 + 9\,900 = 40\,470$, arrondie à
\textbf{40\,000\keur}.

\noindent $\mathrm{PT} = \mathrm{BE} + \mathrm{MR} = 920\,000 + 40\,000 = \mathbf{960\,000}$\keur,
soit $+10\,000$\keur{} par rapport aux comptes sociaux : la meilleure estimation est plus basse que
la provision comptable ($-30\,000$), mais la marge de risque fait plus que compenser.
La méthode exacte (règlement délégué, art.~37) actualise année par année le coût du capital
immobilisé sur le SCR projeté jusqu'à extinction des engagements.
\end{corrige}

\begin{exercice}[Impôts différés]
Les écarts de valorisation sont des différences temporaires avec la base fiscale. Calculer
l'impôt différé correspondant et dire s'il figure à l'actif ou au passif du bilan prudentiel.
\end{exercice}
\begin{corrige}
Actifs plus bas ($-26\,000$) et provisions plus hautes ($+10\,000$) qu'en base fiscale : 36\,000\keur
de différences déductibles, donc un \textbf{impôt différé actif} de
$25\,\% \times 36\,000 = \mathbf{9\,000}$\keur{} (règlement délégué, art.~15). Il n'est retenu que
s'il est recouvrable, hypothèse donnée ici.
\end{corrige}

\begin{exercice}[Bilan prudentiel]
\begin{enumerate}[nosep]
  \item Construire le bilan prudentiel et calculer l'excédent d'actif sur passif.
  \item Le retrouver à partir des capitaux propres comptables.
\end{enumerate}
\end{exercice}
\begin{corrige}
\begin{center}
\begin{tabular}{lr@{\qquad}lr}
\toprule
Actif & & Passif & \\
\midrule
Placements et trésorerie (ex.~1.1) & 1\,025\,000 & Meilleure estimation & 920\,000 \\
Créances                           &    19\,000 & Marge de risque      &  40\,000 \\
Impôts différés actifs             &     9\,000 & Dettes subordonnées  &  50\,000 \\
                                   &            & Autres dettes        &  20\,000 \\
                                   &            & \textbf{Excédent d'actif sur passif} & \textbf{23\,000} \\
\midrule
Total                              & 1\,053\,000 & Total & 1\,053\,000 \\
\bottomrule
\end{tabular}
\end{center}
Réconciliation : $50\,000$ (capitaux propres comptables) $- 26\,000$ (actifs) $- 10\,000$
(provisions) $+ 9\,000$ (impôts différés) $= 23\,000$\keur. La dette subordonnée reste au passif :
elle n'entre pas dans l'excédent, mais dans les fonds propres de base (exercice suivant).
\end{corrige}

\begin{exercice}[Fonds propres]
\begin{enumerate}[nosep]
  \item Calculer les fonds propres de base (directive, art.~88).
  \item Les répartir en niveaux 1, 2 et 3. Calculer la réserve de réconciliation.
\end{enumerate}
\end{exercice}
\begin{corrige}
Fonds propres de base = excédent + passifs subordonnés $= 23\,000 + 50\,000 = \mathbf{73\,000}$\keur.

\noindent Niveau 3 : impôts différés actifs nets, 9\,000. Niveau 2 : dette subordonnée, 50\,000.
Niveau 1 : capital social (30\,000) et réserve de réconciliation, qui est l'excédent diminué des
autres éléments de fonds propres (règlement délégué, art.~70) :
$23\,000 - 30\,000 - 9\,000 = \mathbf{-16\,000}$\keur. D'où
\textbf{niveau~1 = 14\,000}\keur. Une réserve de réconciliation négative signifie que les
ajustements prudentiels consomment une partie du capital social.
\end{corrige}

\begin{exercice}[Couverture du SCR et du MCR]
\begin{enumerate}[nosep]
  \item Appliquer les limites d'éligibilité (règlement délégué, art.~82) et calculer les ratios de
        couverture du SCR et du MCR.
  \item Quel apport en capital (niveau 1) rétablirait la couverture du SCR ?
\end{enumerate}
\end{exercice}
\begin{corrige}
\textbf{SCR} (70\,000) : niveau 3 $< 15\,\%$ du SCR (10\,500) : 9\,000 éligibles ; niveaux 2 + 3
$\leq 50\,\%$ (35\,000) : niveau 2 éligible $35\,000 - 9\,000 = 26\,000$ ; niveau 1 : 14\,000.
Fonds propres éligibles $= 49\,000$, \textbf{ratio $= 70\,\%$}. De plus, le niveau 1 devrait
représenter au moins la moitié du SCR (35\,000) : la couverture échoue aussi sur la qualité.

\noindent\textbf{MCR} (17\,500) : niveau 1 $\geq 80\,\%$ du MCR (14\,000) : tout juste respecté ;
niveau 2 $\leq 20\,\%$ (3\,500) ; niveau 3 exclu. Éligibles $= 17\,500$,
\textbf{ratio $= 100\,\%$}, sans aucune marge.

\noindent\textbf{Apport} : il faut porter le niveau 1 à 35\,000, soit \textbf{21\,000}\keur. Le
niveau 1 atteint alors la moitié du SCR et les éligibles $35\,000 + 35\,000 = 70\,000$ couvrent
le SCR à 100\,\%. L'entreprise doit présenter un plan de rétablissement (directive, art.~138).
\end{corrige}

\begin{exercice}[Sensibilité à trois chocs]
SCR et MCR inchangés, même traitement fiscal. Pour chaque scénario, calculer l'effet sur
l'excédent d'actif sur passif après impôts différés, puis le nouveau niveau 1.
\begin{enumerate}[nosep]
  \item[A.] Actions $-20\,\%$.
  \item[B.] Taux $+100$\,pb : obligations $-5\,\%$, meilleure estimation $-45\,750$\keur.
  \item[C.] Écarts de crédit $+100$\,pb : obligations d'entreprises $-4\,\%$.
\end{enumerate}
Quel scénario est le plus pénalisant ? Le MCR reste-t-il couvert ? Pourquoi la hausse des taux
améliore-t-elle ici la situation ?
\end{exercice}
\begin{corrige}
\begin{center}
\begin{tabular}{lrrrrr}
\toprule
Scénario & Avant impôt & Après impôt & Excédent & IDA nets & Niveau 1 \\
\midrule
A. Actions $-20\,\%$   & $-18\,000$ & $-13\,500$ &  9\,500 & 13\,500 & $-4\,000$ \\
B. Taux $+100$\,pb      & $+12\,000$ &  $+9\,000$ & 32\,000 &  6\,000 & 26\,000 \\
C. Spread $+100$\,pb    &  $-7\,800$ &  $-5\,850$ & 17\,150 & 10\,950 &  6\,200 \\
\bottomrule
\end{tabular}
\end{center}
Le choc actions est le plus pénalisant : le niveau 1 devient négatif, le MCR n'est plus couvert
(il faudrait 14\,000 de niveau 1). L'entreprise doit soumettre un plan de financement à court
terme (directive, art.~139) et s'expose au retrait d'agrément. En B, les provisions baissent plus
que les obligations : la duration du passif dépasse celle de l'actif, et ce même écart rendrait
une baisse des taux dangereuse. C'est l'objet de l'adossement actif-passif (séances 5 et 6).
\end{corrige}

\begin{generalisation}{Ce que le TD simplifie}
Le SCR est donné et ne bouge pas après les chocs ; aucune absorption des pertes par la
participation aux bénéfices ; la base fiscale est celle des comptes sociaux ; la dette subordonnée
garde sa valeur. Les séances 4 à 7 lèvent ces simplifications.
\end{generalisation}

\end{document}
